% !TEX spellcheck = en_US



% ~~~~~~~~~~~~~~~~~~~~~~~~~~~~~~~~~~~~~~~ V
% (NC) 2026 Haluk Bingol
% github.com/halukbingol/LaTeX-Templates
%
% Licensees may copy, distribute, display, and perform the work and make derivative works 
% and remixes based on it only for non-commercial purposes. 
% ~~~~~~~~~~~~~~~~~~~~~~~~~~~~~~~~~~~~~~~ A




% =====================================================================
%  The Eclipse IDE -- Junior Level: Compiler Settings, Null Analysis, Builds and Remote Debugging
%  ---------------------------------------------------------------------
%  Built on hblib/Template/hblibTemplate.tex with the libraries in ../../hblib/.
%  Programs and their real outputs are in codeoutput/:
%     codeoutput/<Name>.java  a program (classes without main() are only compiled)
%     codeoutput/<Name>.sh    a terminal session, run in a copy of codeoutput/<Name>/
%     codeoutput/<Name>.txt   what it printed, made by:  bash run_examples.sh
% =====================================================================

\documentclass[aspectratio=169,10pt,compress]{beamer}

% ~~~~~~~~~~~~~~~~~~~~~~~~~~~~~~~~~~~~~~~~~~~~~~~~~~~~~~~~~~~~~~~~~ V hblib
	\usepackage{../../hblib/hblibPresentation} % lib presentation: theme, i/N, \hSection, algorithm2e, ...
	\usepackage{../../hblib/hblibCodeoutput} % lib code + output: \lstcode, \lstterminal, ...
	\usepackage{../../hblib/hblibDatastructures} % lib data structures: \hString, \hStack, \hQueue
% ~~~~~~~~~~~~~~~~~~~~~~~~~~~~~~~~~~~~~~~~~~~~~~~~~~~~~~~~~~~~~~~~~ A hblib


% theme, i/N, fonts (T1, lmodern), booktabs, xcolor, array (L{width} columns), listings, tikz,
% algorithm2e (\small captions, \Fn) and the color gitred come from the hblib libraries




% +++++++++++++++++++++++++++++++++++++++ V specific
\DontPrintSemicolon % in algorithms
\newcommand{\win}{\textbf{Windows}}
\newcommand{\mac}{\textbf{macOS}}
% +++++++++++++++++++++++++++++++++++++++ A specific

% ---------------------------------------------------------------------
\title{The Eclipse IDE}
\subtitle{Junior Level: Compiler Settings, Null Analysis, Builds and Remote Debugging}
\author[Bingol]{Haluk O. Bingol}

\institute[FCIS]{
    Faculty of Computer and Information Sciences\\
    Yeditepe University
}
\date{
	{\tiny \hbTimeStamp}
}

\begin{document}




% ~~~~~~~~~~~~~~~~~~~~~~~~~~~~~~~~~~~~~~ V frm
\begin{frame}[t,fragile]
  \titlepage
\end{frame}
% ~~~~~~~~~~~~~~~~~~~~~~~~~~~~~~~~~~~~~~ A frm




% ~~~~~~~~~~~~~~~~~~~~~~~~~~~~~~~~~~~~~~ V frm
\begin{frame}[t,fragile] \frametitle{Outline}
  \begin{columns}[T]
    \begin{column}{0.48\textwidth}
      \tableofcontents[sections={1-4}, hideallsubsections]
    \end{column}
    \begin{column}{0.48\textwidth}
      \tableofcontents[sections={5-8}, hideallsubsections]
    \end{column}
  \end{columns}
\end{frame}
% ~~~~~~~~~~~~~~~~~~~~~~~~~~~~~~~~~~~~~~ A frm

% ~~~~~~~~~~~~~~~~~~~~~~~~~~~~~~~~~~~~~~ V
\hSection[Compiler]{Compiler Settings}
% ~~~~~~~~~~~~~~~~~~~~~~~~~~~~~~~~~~~~~~ A




% ~~~~~~~~~~~~~~~~~~~~~~~~~~~~~~~~~~~~~~ V frm
\begin{frame}[t,fragile] \frametitle{Errors, warnings, ignore: you decide}
  \begin{columns}[T]
    \begin{column}{0.5\textwidth}
      \begin{itemize}\small
        \item \menu{Window > Preferences > Java > Compiler > Errors/Warnings} (\mac: \menu{Eclipse > Settings\dots}):
              many checks, each \emph{Error}, \emph{Warning}, \emph{Info} or \emph{Ignore}.
        \item These are the workspace settings: they are \textbf{yours}, not the project's.
        \item Better for a team: project properties, \menu{Java Compiler > Errors/Warnings}, tick
              \emph{Enable project specific settings}.
        \item Eclipse writes them to \texttt{.settings/org.eclipse.jdt.core.prefs}: commit it, and everybody gets the same checks.
      \end{itemize}
    \end{column}
    \begin{column}{0.46\textwidth}
      \lstcode[basicstyle=\ttfamily\tiny]{codeoutput/Defaults/src/Orders.java}{Java}
    \end{column}
  \end{columns}
\end{frame}
% ~~~~~~~~~~~~~~~~~~~~~~~~~~~~~~~~~~~~~~ A frm




% ~~~~~~~~~~~~~~~~~~~~~~~~~~~~~~~~~~~~~~ V frm
\begin{frame}[t,fragile] \frametitle{The defaults}
  \lstterminal[basicstyle=\ttfamily\tiny]{codeoutput/Defaults.txt}

  \vspace{0.1em}
  \small
  With the default settings: two raw types and two unused locals are warnings; the possible \texttt{null} in line 12 is \textbf{not} reported.
\end{frame}
% ~~~~~~~~~~~~~~~~~~~~~~~~~~~~~~~~~~~~~~ A frm




% ~~~~~~~~~~~~~~~~~~~~~~~~~~~~~~~~~~~~~~ V frm
\begin{frame}[t,fragile] \frametitle{The same code with the project's settings file}
  \begin{columns}[T]
    \begin{column}{0.49\textwidth}
      \lstterminal[basicstyle=\ttfamily\tiny, lastline=17]{codeoutput/Prefs.txt}
    \end{column}
    \begin{column}{0.49\textwidth}
      \lstterminal[basicstyle=\ttfamily\tiny, firstline=18]{codeoutput/Prefs.txt}
    \end{column}
  \end{columns}

  \vspace{0.1em}
  \small
  \texttt{ecj -properties} reads the file Eclipse writes, so the command line, a build server and every team member's
  Eclipse report the same problems.
\end{frame}
% ~~~~~~~~~~~~~~~~~~~~~~~~~~~~~~~~~~~~~~ A frm

% ~~~~~~~~~~~~~~~~~~~~~~~~~~~~~~~~~~~~~~ V
\hSection[Null]{Null Analysis}
% ~~~~~~~~~~~~~~~~~~~~~~~~~~~~~~~~~~~~~~ A




% ~~~~~~~~~~~~~~~~~~~~~~~~~~~~~~~~~~~~~~ V frm
\begin{frame}[t,fragile] \frametitle{Contracts with \texttt{@NonNull} and \texttt{@Nullable}}
  \begin{columns}[T]
    \begin{column}{0.44\textwidth}
      \lstcode[basicstyle=\ttfamily\tiny]{codeoutput/NullAnn/src/app/Names.java}{Java}
    \end{column}
    \begin{column}{0.54\textwidth}
      \uncover<2->{\lstterminal[basicstyle=\ttfamily\tiny]{codeoutput/NullAnn.txt}}
    \end{column}
  \end{columns}
  \small \texttt{@NonNull} parameter: callers must not pass \texttt{null}. \texttt{@Nullable} result: callers must check it first.
\end{frame}
% ~~~~~~~~~~~~~~~~~~~~~~~~~~~~~~~~~~~~~~ A frm




% ~~~~~~~~~~~~~~~~~~~~~~~~~~~~~~~~~~~~~~ V frm
\begin{frame}[t,fragile] \frametitle{The annotations and the settings}
  \begin{columns}[T]
    \begin{column}{0.48\textwidth}
      \lstcode[basicstyle=\ttfamily\tiny]{codeoutput/NullAnn/src/ann/NonNull.java}{Java}
    \end{column}
    \begin{column}{0.48\textwidth}
      \lstcode[basicstyle=\ttfamily\tiny]{codeoutput/NullAnn/null.prefs}{}
      \begin{itemize}\small
        \item \emph{Errors/Warnings $\to$ Null analysis $\to$ Enable annotation-based null analysis}.
        \item Eclipse ships its own annotations (\texttt{org.eclipse.jdt.annotation}); Quick Fix adds the JAR.
              Other names (e.g.\ JSpecify's) can be configured; here: two tiny annotations of our own.
        \item The analysis is at \textbf{compile time}: no cost when the program runs.
      \end{itemize}
    \end{column}
  \end{columns}
\end{frame}
% ~~~~~~~~~~~~~~~~~~~~~~~~~~~~~~~~~~~~~~ A frm

% ~~~~~~~~~~~~~~~~~~~~~~~~~~~~~~~~~~~~~~ V
\hSection[Builds]{Maven and Gradle Projects}
% ~~~~~~~~~~~~~~~~~~~~~~~~~~~~~~~~~~~~~~ A




% ~~~~~~~~~~~~~~~~~~~~~~~~~~~~~~~~~~~~~~ V frm
\begin{frame}[t,fragile] \frametitle{Eclipse and build tools}
  \begin{columns}[T]
    \begin{column}{0.5\textwidth}


% +++++++++++++++++++++++++++++++++++++++ V lst
%: -Lst.
\begin{lstlisting}[style=codesnippet, language=XML, basicstyle=\ttfamily\tiny]
<project>
  <modelVersion>4.0.0</modelVersion>
  <groupId>tr.edu.yeditepe</groupId>
  <artifactId>grades</artifactId>
  <version>1.0</version>
  <properties>
    <maven.compiler.release>21</maven.compiler.release>
  </properties>
  <dependencies>
    <dependency>
      <groupId>org.junit.jupiter</groupId>
      <artifactId>junit-jupiter</artifactId>
      <version>5.12.2</version>
      <scope>test</scope>
    </dependency>
  </dependencies>
</project>
\end{lstlisting}
% +++++++++++++++++++++++++++++++++++++++ A lst


    \end{column}
    \begin{column}{0.46\textwidth}
      \begin{itemize}\small
        \item \textbf{m2e} (Maven) and \textbf{Buildship} (Gradle) come with the Java package.
        \item \menu{File > Import > Maven > Existing Maven Projects}: Eclipse reads \texttt{pom.xml},
              downloads the libraries and \textbf{writes} the build path for you.
        \item After editing \texttt{pom.xml}: \menu{Maven > Update Project} (\keys{Alt+F5} / \keys{Option+F5}).
        \item The build tool is the truth: the build server runs \hCode{mvn verify}, not Eclipse.
      \end{itemize}
      \small \hRemark{Not run here}: Maven needs the internet to download the libraries.
    \end{column}
  \end{columns}
\end{frame}
% ~~~~~~~~~~~~~~~~~~~~~~~~~~~~~~~~~~~~~~ A frm

% ~~~~~~~~~~~~~~~~~~~~~~~~~~~~~~~~~~~~~~ V
\hSection[Errors]{Code With Errors Still Runs}
% ~~~~~~~~~~~~~~~~~~~~~~~~~~~~~~~~~~~~~~ A




% ~~~~~~~~~~~~~~~~~~~~~~~~~~~~~~~~~~~~~~ V frm
\begin{frame}[t,fragile] \frametitle{``Errors exist \dots{} Proceed with launch?''}
  \begin{columns}[T]
    \begin{column}{0.46\textwidth}
      \lstcode[basicstyle=\ttfamily\tiny]{codeoutput/Proceed/src/Shop.java}{Java}
    \end{column}
    \begin{column}{0.52\textwidth}
      \uncover<2->{\lstterminal[basicstyle=\ttfamily\tiny]{codeoutput/Proceed.txt}}
    \end{column}
  \end{columns}
  \small Unlike \texttt{javac}, \texttt{ecj} still writes the class: the broken method throws \texttt{Unresolved compilation problem}
  when it is \textbf{called}. Eclipse uses this; that is why \emph{Proceed} ``works'' until it does not.
\end{frame}
% ~~~~~~~~~~~~~~~~~~~~~~~~~~~~~~~~~~~~~~ A frm

% ~~~~~~~~~~~~~~~~~~~~~~~~~~~~~~~~~~~~~~ V
\hSection[Compilers]{ecj and javac}
% ~~~~~~~~~~~~~~~~~~~~~~~~~~~~~~~~~~~~~~ A




% ~~~~~~~~~~~~~~~~~~~~~~~~~~~~~~~~~~~~~~ V frm
\begin{frame}[t,fragile] \frametitle{Two compilers, one language: javac}
  \begin{columns}[T]
    \begin{column}{0.4\textwidth}
      \lstcode[basicstyle=\ttfamily\tiny]{codeoutput/Compare/Legacy.java}{Java}
      \begin{itemize}\small
        \item Both follow the Java Language Specification; the class files run on the same JVM.
        \item Different programs: messages, warnings and some corner cases differ.
      \end{itemize}
    \end{column}
    \begin{column}{0.58\textwidth}
      \uncover<2->{\lstterminal[basicstyle=\ttfamily\tiny, lastline=11]{codeoutput/Compare.txt}}
    \end{column}
  \end{columns}
\end{frame}
% ~~~~~~~~~~~~~~~~~~~~~~~~~~~~~~~~~~~~~~ A frm




% ~~~~~~~~~~~~~~~~~~~~~~~~~~~~~~~~~~~~~~ V frm
\begin{frame}[t,fragile] \frametitle{Two compilers, one language: ecj}
  \begin{columns}[T]
    \begin{column}{0.4\textwidth}
      \lstcode[basicstyle=\ttfamily\tiny]{codeoutput/Compare/Legacy.java}{Java}
      \begin{itemize}\small
        \item Both follow the Java Language Specification; the class files run on the same JVM.
        \item Different programs: messages, warnings and some corner cases differ.
      \end{itemize}
    \end{column}
    \begin{column}{0.58\textwidth}
      \uncover<2->{\lstterminal[basicstyle=\ttfamily\tiny, firstline=12]{codeoutput/Compare.txt}}
    \end{column}
  \end{columns}
\end{frame}
% ~~~~~~~~~~~~~~~~~~~~~~~~~~~~~~~~~~~~~~ A frm

% ~~~~~~~~~~~~~~~~~~~~~~~~~~~~~~~~~~~~~~ V
\hSection[Debugging]{Advanced Debugging}
% ~~~~~~~~~~~~~~~~~~~~~~~~~~~~~~~~~~~~~~ A




% ~~~~~~~~~~~~~~~~~~~~~~~~~~~~~~~~~~~~~~ V frm
\begin{frame}[t,fragile] \frametitle{Remote debugging}
  \begin{columns}[T]
    \begin{column}{0.5\textwidth}
      \begin{itemize}\small
        \item Start the program with the debug agent: it \textbf{listens} on a port (here 5005) and waits (\texttt{suspend=y}).
        \item Eclipse: \menu{Run > Debug Configurations\dots}, \emph{Remote Java Application}, host and port 5005, \textbf{Debug}.
        \item Works for a server, a container, another computer \dots{} \hRemark{Never open a debug port to the internet: whoever connects controls the program.}
      \end{itemize}
    \end{column}
    \begin{column}{0.46\textwidth}
      \begin{center}\footnotesize
      \begin{tabular}{@{}L{0.46\linewidth}L{0.5\linewidth}@{}}
        \toprule
        \textbf{agent option}  & \textbf{meaning} \\
        \midrule
        \texttt{transport=dt\_socket} & TCP connection \\
        \texttt{server=y}      & the program waits for the debugger \\
        \texttt{suspend=y}     & do not start before it connects \\
        \texttt{address=5005}  & the port (\texttt{localhost} only by default) \\
        \bottomrule
      \end{tabular}
      \end{center}
    \end{column}
  \end{columns}
\end{frame}
% ~~~~~~~~~~~~~~~~~~~~~~~~~~~~~~~~~~~~~~ A frm




% ~~~~~~~~~~~~~~~~~~~~~~~~~~~~~~~~~~~~~~ V frm
\begin{frame}[t,fragile] \frametitle{Attaching to a waiting program}
  \begin{columns}[T]
    \begin{column}{0.49\textwidth}
      \lstterminal[basicstyle=\ttfamily\tiny, lastline=16]{codeoutput/Remote.txt}
    \end{column}
    \begin{column}{0.49\textwidth}
      \lstterminal[basicstyle=\ttfamily\tiny, firstline=17]{codeoutput/Remote.txt}
    \end{column}
  \end{columns}

  \vspace{0.1em}
  \small
  \texttt{jdb -attach} does what Eclipse's \emph{Remote Java Application} does. The program printed its own lines into \texttt{server.log}.
\end{frame}
% ~~~~~~~~~~~~~~~~~~~~~~~~~~~~~~~~~~~~~~ A frm




% ~~~~~~~~~~~~~~~~~~~~~~~~~~~~~~~~~~~~~~ V frm
\begin{frame}[t,fragile] \frametitle{More debugger tools}
  \begin{center}\small
  \begin{tabular}{L{0.3\textwidth}L{0.62\textwidth}}
    \toprule
    \textbf{tool}                       & \textbf{use} \\
    \midrule
    Exception breakpoint                & stop where an exception is thrown, even if it is caught later \\
    Watchpoint (on a field)             & stop whenever a field is read or changed: who sets it to \texttt{null}? \\
    Method breakpoint                   & stop on entry to or exit from a method (slower) \\
    \emph{Drop to Frame}                & restart the current method with the current values \\
    Hot code replace                    & edit a method body while stopped, save, continue with the new code \\
    \emph{Display} view                 & run any Java snippet in the stopped program \\
    \bottomrule
  \end{tabular}
  \end{center}
  \small Hot code replace works only for method bodies; a change of fields or signatures needs a restart.
\end{frame}
% ~~~~~~~~~~~~~~~~~~~~~~~~~~~~~~~~~~~~~~ A frm

% ~~~~~~~~~~~~~~~~~~~~~~~~~~~~~~~~~~~~~~ V
\hSection[Navigation]{Working in Large Code}
% ~~~~~~~~~~~~~~~~~~~~~~~~~~~~~~~~~~~~~~ A




% ~~~~~~~~~~~~~~~~~~~~~~~~~~~~~~~~~~~~~~ V frm
\begin{frame}[t,fragile] \frametitle{Finding your way}
  \begin{center}\small
  \begin{tabular}{llll}
    \toprule
    \textbf{question}                    & \textbf{menu}                                & \win                  & \mac \\
    \midrule
    where is class \dots?                & \menu{Navigate > Open Type\dots}             & \keys{Ctrl+\shift+T}  & \keys{\cmd+\shift+T} \\
    where is file \dots?                 & \menu{Navigate > Open Resource\dots}         & \keys{Ctrl+\shift+R}  & \keys{\cmd+\shift+R} \\
    what is in this class?               & \menu{Navigate > Quick Outline}              & \keys{Ctrl+O}         & \keys{\cmd+O} \\
    who uses this method?                & \menu{Search > References > Workspace}       & \keys{Ctrl+\shift+G}  & \keys{\cmd+\shift+G} \\
    who calls whom?                      & \menu{Navigate > Open Call Hierarchy}        & \keys{Ctrl+Alt+H}     & menu \\
    sub- and superclasses?               & \menu{Navigate > Open Type Hierarchy}        & \keys{F4}             & \keys{F4} \\
    text anywhere?                       & \menu{Search > File\dots}                    & \keys{Ctrl+H}         & \keys{Ctrl+H} \\
    \bottomrule
  \end{tabular}
  \end{center}
  \small \textbf{Working sets} (Package Explorer menu, \emph{Select Working Set\dots}) hide the projects you are not working on.
\end{frame}
% ~~~~~~~~~~~~~~~~~~~~~~~~~~~~~~~~~~~~~~ A frm




% ~~~~~~~~~~~~~~~~~~~~~~~~~~~~~~~~~~~~~~ V frm
\begin{frame}[t,fragile] \frametitle{When Eclipse itself misbehaves}
  \begin{columns}[T]
    \begin{column}{0.48\textwidth}
      \begin{block}{\texttt{eclipse.ini}: memory and Java}
        \small
        \win: next to \texttt{eclipse.exe}.\\
        \mac: \texttt{Eclipse.app/Contents/Eclipse/eclipse.ini}.\\
        After the line \texttt{-vmargs}, e.g.\ \texttt{-Xmx4g} for big workspaces.
      \end{block}
      \begin{block}{The log}
        \small \texttt{<workspace>/.metadata/.log}, or \menu{Help > About > Installation Details > Configuration > View Error Log}.
      \end{block}
    \end{column}
    \begin{column}{0.48\textwidth}
      \begin{block}{In this order}
        \small
        \begin{enumerate}
          \item \menu{Project > Clean\dots}: rebuild everything.
          \item \menu{File > Refresh} (\keys{F5}): files changed outside Eclipse.
          \item \menu{Maven > Update Project}.
          \item Restart with \hCode{-clean} (plug-in caches).
          \item New workspace, \menu{File > Import > Existing Projects into Workspace}.
        \end{enumerate}
      \end{block}
    \end{column}
  \end{columns}
\end{frame}
% ~~~~~~~~~~~~~~~~~~~~~~~~~~~~~~~~~~~~~~ A frm

% ~~~~~~~~~~~~~~~~~~~~~~~~~~~~~~~~~~~~~~ V
\hSection[Summary]{Summary}
% ~~~~~~~~~~~~~~~~~~~~~~~~~~~~~~~~~~~~~~ A




% ~~~~~~~~~~~~~~~~~~~~~~~~~~~~~~~~~~~~~~ V frm
\begin{frame}[t,fragile] \frametitle{Recipes}
  \begin{center}\footnotesize
  \begin{tabular}{@{}L{0.36\textwidth}L{0.58\textwidth}@{}}
    \toprule
    \textbf{goal}                              & \textbf{recipe} \\
    \midrule
    same checks for the whole team             & project-specific compiler settings; commit \texttt{.settings/} \\
    no more \texttt{NullPointerException}s     & annotation-based null analysis, errors not warnings \\
    libraries without manual build paths       & Maven or Gradle project, \menu{Maven > Update Project} \\
    debug a program that runs elsewhere        & \texttt{-agentlib:jdwp=\dots,server=y}, \emph{Remote Java Application} \\
    find a bug that appears after an hour      & conditional breakpoint, watchpoint, exception breakpoint \\
    Eclipse acts strange                       & Clean, Refresh, Update Project, \texttt{-clean}, the \texttt{.log} \\
    \bottomrule
  \end{tabular}
  \end{center}
\end{frame}
% ~~~~~~~~~~~~~~~~~~~~~~~~~~~~~~~~~~~~~~ A frm




% ~~~~~~~~~~~~~~~~~~~~~~~~~~~~~~~~~~~~~~ V frm
\begin{frame}[t,fragile] \frametitle{Exercises}
  \begin{enumerate}\small
    \item Make ``potential null pointer access'' an error in a project of yours; commit \texttt{.settings/} and let a friend import the project.
    \item Annotate a class with \texttt{@NonNull}/\texttt{@Nullable} (Eclipse's or JSpecify's) until there are no null warnings left.
    \item Import a Maven project from GitHub into Eclipse; find where the downloaded JARs are stored.
    \item Start a program with the debug agent in a terminal and attach Eclipse's debugger to it.
    \item Use a watchpoint to find which method sets a field to \texttt{null} in a program of 5 or more classes.
  \end{enumerate}
\end{frame}
% ~~~~~~~~~~~~~~~~~~~~~~~~~~~~~~~~~~~~~~ A frm




% ~~~~~~~~~~~~~~~~~~~~~~~~~~~~~~~~~~~~~~ V frm
\begin{frame}[t,fragile]
	\begin{center}
		\vspace{4cm}
		\hRemark{\Huge Questions?}
	\end{center}
\end{frame}
% ~~~~~~~~~~~~~~~~~~~~~~~~~~~~~~~~~~~~~~ A frm



\end{document}
